%!TEX root = ./../../springer-mono-format/book.tex
%!TEX root = ./../../thesis-trier.tex

To every (complete and deterministic) semiautomaton $\semiautomat{A}{Q}{\Sigma}{\delta}$
corresponds its transformation semigroup $\TS{\mathcal A}$
with generating set $\{ \fword{\delta}{a} \colon Q \to Q \mid a \in \Sigma \}$.
Conversely, let $S \le_S \mathcal T_n$ 
be a transformation semigroup
generated by $g_1, \ldots, g_m$. Then for  $\Sigma = \{g_1, \ldots, g_n\}$ (viewing the generating sequence as a set of abstract symbols)
we obtain
a semiautomaton $\semiautomat{A}{[n]}{\Sigma}{\delta}$
with $\TS{\mathcal A} = S$
by setting $\applytrans{\delta}{q}{g_i} = \applyfunc{q}{g_i}$
for $q \in [n]$ and $i \in \{1,\ldots,m\}$.\footnote{A transformation $g_i$ is allowed to appear more than once among  $g_1,\ldots,g_m$.} 
Thus, a semiautomaton may be viewed as a transformation semigroup together
with a distinguished generating sequence.

\begin{example}	\label{prelim:aut:ex:Z6genset}
	The choice of a generating set for a given transformation semigroup $S \le_S \mathcal T_n$
	can lead to quite distinct automata.
	For example, consider the permutation group $G = \langle g \rangle_G$ with $g = (0 ~ 1 ~ \cdots ~ 5)$.
	When constructing the automaton with generating set $\{ g \}$ 
	we obtain an automaton that cyclically permutes its states.
	Another generating set is given by $g^2 = (0 ~ 2 ~ 4) (1 ~ 3 ~ 5)$ and $g^3 = (0 ~ 3) (1 ~ 4) (2 ~ 5)$
	since $g^2 g^2 g^3 = g$.
	The automaton we obtain from this generating set is depicted in \cref{prelim:aut:fig:Z6genset}.	
	
	\begin{figure}[ht]
		\centering	
		\begin{tikzpicture}[node distance = 6em and 10em, on grid, auto, ->]
			\node (q0) [state] {$0$};
			\node (q2) [state, right = of q0] {$2$};
			\node (q4) [state, right = of q2] {$4$};
			\node (q3) [state, below = of q0] {$3$};
			\node (q5) [state, below = of q2] {$5$};
			\node (q1) [state, below = of q4] {$1$};
			
			\path (q0) edge node {$g^2$} (q2)
			(q2) edge node {$g^2$} (q4)
			(q4) edge [bend right=23,above] node {$g^2$} (q0);
			
			\path (q3) edge node {$g^2$} (q5)
			(q5) edge node {$g^2$} (q1)
			(q1) edge [bend left=23] node {$g^2$} (q3);			
			
			\path (q0) edge [bend left] node {$g^3$} (q3)
			(q3) edge [bend left] node {$g^3$} (q0);
			
			\path (q2) edge [bend left] node {$g^3$} (q5)
			(q5) edge [bend left] node {$g^3$} (q2);
						
			\path (q4) edge [bend left] node {$g^3$} (q1)
			(q1) edge [bend left] node {$g^3$} (q4);
			
		\end{tikzpicture}
		\caption[Automaton constructed from transformation semigroup.]{Automaton constructed from permutation group, see \cref{prelim:aut:ex:Z6genset}.}
		\label{prelim:aut:fig:Z6genset}
	\end{figure}		
\end{example}

Many of our notions and results
can be formulated either for automata or for transformation semigroups.
For example, $\mathcal A$ is synchronizing if and only if $\TS{\mathcal A}$
contains a constant transformation.

Our results are motivated by automata-theoretical problems, but properties only
related to subset mapping properties do not depend on the choice of some generating set 
and are actually properties of transformation
semigroups. So, these properties are better expressed in an algebraic setting
with transformation semigroups, as done in \cref{part:aut}. The corresponding automata-theoretic 
reformulations are usually straightforward and are therefore omitted.

A semiautomaton is a machine model that processes, or reacts to, inputs.
In this setting, the alphabet is crucial, as it is the means by which
the machine interacts with its environment.
Since inputs correspond to physical actions and an operator may want
to minimize the effort to achieve a given goal,
it is natural to ask for the minimum number of such 
actions required to do this.
Thus, subset mapping properties with length bounds on shortest words
realizing some property are naturally expressed in an automata-theoretic formulation.

Later in \cref{part:constrsync}, we use automata-theoretic notation exclusively.

We note that there is a generalization of the \Cerny conjecture for transformation monoids~\cite{AnanichevVolkov04,Rystsov2024}.
Shortest words in some generating sets have also been investigated
in the context of the diameter of a Cayley graph depends~\cite{DBLP:journals/ijfcs/GonzeGJGV19}.
%Coxeter groups are intimately related to their generating sets and the element represented
%by the longest possible word in the generators~\cite{}, 
%or the set of words representing geodesics in geometric group theory~\cite{jamesworringten}.